\documentclass[12pt,a4paper]{article}
\usepackage[russian]{babel}
\usepackage{fontspec}
\usepackage{geometry}
\usepackage{booktabs}
\usepackage{longtable}
\usepackage{array}
\usepackage{multirow}
\usepackage{enumitem}
\usepackage{titlesec}
\usepackage{xcolor}
\usepackage{siunitx}

\usepackage{pdfpages}

\usepackage{tablefootnote}

% Дополнительные пакеты для данного документа
\usepackage{caption}
\usepackage{subcaption}
\usepackage{listings}
% \usepackage{hyperref}

\geometry{a4paper, margin=25mm}
\setmainfont{Liberation Serif}
\setsansfont{Liberation Sans}

\definecolor{adblue}{RGB}{0,84,166}

\titleformat{\section}{\large\bfseries\color{adblue}}{}{0em}{}[\titlerule]
\titleformat{\subsection}{\normalsize\bfseries}{}{0em}{}
\titleformat{\subsubsection}{\normalsize\itshape}{}{0em}{}

\sisetup{output-decimal-marker={,}}

\usepackage{tikz}
\usetikzlibrary{positioning, calc, backgrounds, math, 3d, perspective, arrows.meta, graphs, graphdrawing, fit, shapes.geometric, trees}
    

\usepackage{pgfplots} % для построения графиков
\pgfplotsset{compat=1.18} % версия для совместимости
\usepackage{tkz-euclide}

% --- Подключение пакета алгоритмов (без флага algo2e, чтобы вернуть \begin{algorithm}) ---
\usepackage[ruled,vlined]{algorithm2e}
%\usepackage{caption}
\newcommand{\Gray}[1]{\textcolor{gray}{#1}} 
% Настройка подписей
\DeclareCaptionLabelFormat{gost}{#1 #2}
\captionsetup[table]{labelformat=gost,labelsep=endash,justification=justified,singlelinecheck=false,font=normal}
\captionsetup[figure]{labelformat=gost,labelsep=endash,justification=centering,font=normal}

% Локализация algorithm2e на русский (ОБЯЗАТЕЛЬНО после загрузки пакета)
\SetAlgorithmName{Алгоритм}{}{}
\SetAlgoCaptionSeparator{---}
\SetKwInput{Input}{Вход}
\SetKwInput{Output}{Выход}

\usepackage{url}
% Говорим пакету использовать моноширинный шрифт (как у texttt)
\Urlmuskip=0mu plus 1mu % гибкие пробелы для формул, если нужно
\DeclareUrlCommand\code{\urlstyle{tt}} 

% Библиотека схемотехники
\usepackage[siunitx, RPvoltages, american]{circuitikzgit}
\ctikzset{
    amplifiers/fill=cyan,
    sources/fill=green!20,
    diodes/fill=white,
    resistors/fill=violet,
    grounds/scale=1.2,
}

\usepackage{iftex}

\ifXeTeX % Если используется TeTeX или TeLaTeX
  \typeout{Режим XeTeX}
  \usepackage{xltxtra}
  \usepackage{fontspec}
  \usepackage{polyglossia}
  \setmainlanguage{russian}
  \setotherlanguage{english}
  \setkeys{russian}{babelshorthands=true} % По идее включает <<, >>, -- ---

  \setmainfont{Times New Roman}[Ligatures=TeX]
  \setromanfont{Times New Roman}[Ligatures=TeX]
  \setsansfont{Arial}[Ligatures=TeX]
  \setmonofont{Courier New}[Ligatures=TeX] 

  \newfontfamily{\cyrillicfont}{Times New Roman}
  \newfontfamily{\cyrillicfontrm}{Times New Roman}
  \newfontfamily{\cyrillicfonttt}{Courier New}
  \newfontfamily{\cyrillicfontsf}{Arial}
\fi

\ifLuaTeX % Если используется LuaLatex
  \typeout{Режим LuaTeX}
%  \usepackage{luatextra}
  \usepackage{fontspec}
  \defaultfontfeatures{Ligatures={TeX}}
  \usepackage{polyglossia}
  \setmainlanguage{russian}
  \setotherlanguage{english}
  \setmainfont{Times New Roman}[Script=Cyrillic, Ligatures=TeX]
  \setromanfont{Times New Roman}[Script=Cyrillic, Ligatures=TeX]
  \setsansfont{Arial}[Script=Cyrillic, Ligatures=TeX]
  \setmonofont{Courier New}[Ligatures=TeX]
  \newfontfamily{\cyrillicfont}{Times New Roman}[Script=Cyrillic, Ligatures=TeX]
  \newfontfamily{\cyrillicfontrm}{Times New Roman}[Script=Cyrillic, Ligatures=TeX]
  \newfontfamily{\cyrillicfonttt}{Courier New}
  \newfontfamily{\cyrillicfontsf}{Arial}[Script=Cyrillic, Ligatures=TeX]

  \usegdlibrary{trees}
  \usegdlibrary{force}
\fi

\ifPDFTeX % Если используется 8битный PDFTex
  \typeout{8-битный режим pdfTex}
  % For pdfTeX we must use old
  % 8-bit font technologies
  \usepackage[T2A]{fontenc}
  \usepackage[english,russian]{babel}
  \usepackage[utf8]{inputenc}
  \usepackage[english, russian]{babel}

  %\usepackage{newtxmath}
  \usepackage{amsmath}
  \usepackage{amsthm}
  \usepackage{amssymb}
  \usepackage{amsfonts}
  \usepackage{mathtext}
\fi

\usepackage{amsthm}  % Возможности AMSMath
\usepackage{amsmath}
\usepackage{amssymb}

\usepackage{esvect} % Для красивых стрелок
\usepackage{bm} % Для жирных символов

%Hyphenation rules
\usepackage{hyphenat}
\hyphenation{ма-те-ма-ти-ка вос-ста-нав-ли-вать ото-бра-жа-ют-ся сге-не-ри-ро-ван-ном сим-во-лы}
% \usepackage[english, russian]{babel}

\usepackage{graphicx}
\graphicspath{ {./images/} }

\usepackage{empheq}

\usepackage{float}

% \usepackage{comment}

\usepackage[hidelinks,
    unicode=true,          % поддерживает Unicode в закладках
    psdextra=true,      % автоматически конвертирует простую математику в текст
    % focusdefault=false
]{hyperref}


% Окружение "Теорема"
\newtheorem{theorem}{Теорема}[section]
\newtheorem{corollary}{Corollary}[theorem]
\newtheorem{lemma}[theorem]{Lemma}

% Окружение "Определение"
\theoremstyle{definition} % Прямой шрифт, нумерация
\newtheorem{definition}{Определение}[section]
\newtheorem{example}{Пример}[section]

% Окружение "Замечание"
\theoremstyle{remark} % Прямой шрифт, без нумерации
% Окуржение "Решение"
\newtheorem*{solution}{Решение}
\newtheorem*{remark}{Замечание}

\usepackage{polynom}
\usepackage{tabularx}
\usepackage{cancel}

\addto\captionsrussian{%
  \renewcommand{\figurename}{Рис.}%
  \renewcommand{\tablename}{Табл.}%
}

\renewcommand{\arraystretch}{1.5} % Высота строки таблицы

\renewcommand{\labelitemi}{\(\circ\)} % Бюллет -- пустой кружок в itemize

\title{Примечания к предварительной обработке}
\author{Наталия А. Рыбкина}
\date{2026}

\begin{document}

\maketitle

\section{Что делает Байер?}

\subsection{Что такое маска Байера и как она работает?}

Маска Байера — это не «разборка» на три отдельных изображения, а \textbf{одноканальная мозаика}, в которой каждый пиксель хранит значение только одного цвета (R, G или B). Например, для паттерна RGGB:

\[
\begin{array}{|c|c|c|c|}
\hline
R & G_1 & R & G_1 \\
\hline
G_2 & B & G_2 & B \\
\hline
R & G_1 & R & G_1 \\
\hline
G_2 & B & G_2 & B \\
\hline
\end{array}
\]

Зелёный цвет представлен двумя подканалами ($G_1$ и $G_2$), которые в сумме занимают половину площади матрицы. Красный и синий — по четверти. В результате получается \textbf{перемешанный} сигнал, по которому нельзя просто взять и получить три независимых полутоновых изображения. Задача демозаики (восстановления полного RGB) как раз и состоит в том, чтобы по этой мозаике корректно вычислить недостающие цвета в каждой точке.

\section{Ошибки в схеме}

Схема, которую вы привели, отражает \textbf{общую логику}, но содержит несколько \textbf{неточностей} и \textbf{нарушений порядка операций} с точки зрения физической модели, которую мы обосновали ранее. Давайте разберём по пунктам.

\section{Что на схеме указано верно}

- \textbf{Разделение на две ветки} (Подготовка данных → Обучение / Тестирование) – правильно.  
- \textbf{Аугментация RAW-данных} – да, аугментацию нужно делать после формирования мозаики (а не до).  
- \textbf{Сжатие RAW 4x/8x} – как намёк на downscale (уменьшение разрешения) – верно, но важно: сжатие должно применяться к \textbf{полноцветному изображению до маски Байера}, а не к RAW-мозаике.  
- \textbf{Добавление шума} – присутствует, и это хорошо.  

---

\section{Что на схеме неверно или вводит в заблуждение}

\subsection{«RAW-данные» в начале}

На входе у нас должны быть \textbf{эталонные полноцветные изображения высокого качества} (HQ). Это могут быть RAW-файлы, но \textbf{после проявления (demosaicing)} в RGB, либо качественные JPEG/PNG. Называть их «RAW-данные» некорректно, потому что RAW — это уже мозаичный сигнал с сенсора. В эксперименте мы \textbf{синтезируем} этот мозаичный сигнал из полноцветного эталона. Схема должна начинаться с «Эталонное RGB-изображение» или «Full‑colour HQ».

\subsection{«Унификация массива Байера»}

Этот термин мне незнаком. Вероятно, подразумевается приведение разных RAW-форматов к единому паттерну (например, RGGB). Но в нашем конвейере мы \textbf{накладываем маску Байера} на полноцветную картинку. Правильнее назвать этап \textbf{«Наложение маски Байера (RGGB)»}.

\subsection{Порядок операций: сжатие → Байер → шум}

На схеме «Сжатие RAW» стоит \textbf{после} «Унификации массива Байера». Это означает, что вы сначала делаете мозаику, а потом её сжимаете. \textbf{Так делать нельзя}, потому что:
- Сжатие (downscale) должно происходить на эталонном полноцветном изображении \textbf{до} наложения маски, чтобы имитировать физическое уменьшение разрешения (более крупные пиксели сенсора).
- Если сначала сделать мозаику, а потом сжать, то вы нарушите структуру Байера (чередование цветов станет нерегулярным после интерполяции при сжатии).

\textbf{Правильный порядок (как мы обсуждали):}  
\textbf{HQ RGB → Downscale (4x/8x) → PSF blur → Байер → шум → извлечение 4 каналов → аугментация.}

\subsection{Отсутствие этапа «извлечение 4 каналов RGGB»}

После наложения маски Байера и шума у вас получается \textbf{одноканальная мозаика} (2D массив). Нейросеть NAFNet ожидает на входе 4 канала: R, G1, G2, B (отдельно). Этот этап обязателен, на схеме его нет.

\subsection{Ветка «Сжатие RGB → Демозаика RAW - RGB»}

Не понятно, что здесь имеется в виду. Если это этап для тестирования (восстановление из мозаики), то он должен быть после обучения. «Демозаика RAW - RGB» — это как раз то, что делает нейросеть. А «Сжатие RGB» здесь лишнее. Лучше просто: \textbf{Обученная нейросеть → инференс на тестовых мозаиках → сравнение с эталоном (PSNR, SSIM)}.

\subsection{Аугментация должна быть \textbf{после} извлечения 4 каналов}

На схеме «Аугментация RAW-данных» стоит до обучения, но не указано, на каком этапе. Правильно: сформированные 4-канальные патчи (уже в виде тензоров) можно аугментировать (повороты на 90°, отражения) непосредственно в загрузчике данных PyTorch.

---

\section{Исправленная схема (текстовое описание)}

Я нарисую правильную последовательность в виде блоков, которую можно взять за основу для диссертации:

\begin{enumerate}
    \item Эталонное RGB-изображение (HQ)
    \item Пространственное сжатие (downscale 4x или 8x) – имитация мелкого сенсора
    \item Свёртка с PSF (гауссиана) – модель ЧКХ объектива
    \item Наложение маски Байера (паттерн RGGB) – получаем одноканальную мозаику
    \item Добавление коррелированного шума (АБГШ → свёртка с тем же PSF → прибавление)
    \item Извлечение 4-х подканалов (R, G1, G2, B) – формирование 4-канального LQ-патча
    \item Сохранение LQ (16‑бит TIFF) и HQ (PNG после шагов 1–2, без мозаики и шума)
    \item Аугментация (только повороты на 90°, отражения) – применяется на лету при загрузке
    \item Обучение нейросети NAFNet + Focal Frequency Loss
    \item Тестирование на отдельной выборке → расчёт PSNR, SSIM → сравнение с Bilinear, MHC
\end{enumerate}

\section{Вывод по вашей схеме}

\textbf{Схема на картинке является упрощённой и содержит серьёзные ошибки в порядке операций (сжатие после Байера, отсутствие извлечения каналов, неправильное название этапов).}  

Для магистерской диссертации такую схему использовать нельзя. Возьмите за основу исправленную версию, приведённую выше, и при необходимости адаптируйте под свой стиль оформления.


\end{document}